\chapter{Ordres, bons ordres i segments}\label{test:cap13}

Aquest test correspon al capítol~13 de la part teòrica. Cada pregunta té una única resposta correcta.

\iniciTest{u13}{b,c,b,a,c,b,a,c,b,c,b,b}

\begin{enumerate}

\pregunta Un \emph{ordre estricte} sobre \(W\) és una part \({<}\hookrightarrow W\times W\) que és\ldots
\begin{itemize}
\opcio{a}{reflexiva i simètrica.}
\opcio{b}{irreflexiva i transitiva.}
\opcio{c}{una relació d'equivalència.}
\opcio{d}{una aplicació \(W\to W\).}
\end{itemize}
\feedback
\justificacio{Irreflexiva i transitiva: les dues condicions es tallen amb la diagonal i la composició de relacions. La simetria és pròpia de les equivalències, i un ordre no és una aplicació sinó una part de \(W\times W\).}

\pregunta Un ordre estricte és un \emph{bon ordre} quan\ldots
\begin{itemize}
\opcio{a}{és total.}
\opcio{b}{no té elements màxims.}
\opcio{c}{tota part no inicial de \(W\) té un element mínim.}
\opcio{d}{és isomorf a \(\N\).}
\end{itemize}
\feedback
\justificacio{És la definició: tota part no inicial té mínim. La totalitat n'és una conseqüència, no la definició, i hi ha bons ordres molt més llargs que \(\N\).}

\pregunta Tot bon ordre\ldots
\begin{itemize}
\opcio{a}{és finit.}
\opcio{b}{és total (i el mínim d'una part no inicial és únic).}
\opcio{c}{té un element màxim.}
\opcio{d}{admet cadenes descendents infinites.}
\end{itemize}
\feedback
\justificacio{El mínim de \(\{u,w\}\) decideix la comparació sobre punts, i el lema de transferència la converteix en una igualtat de parts; la unicitat ve de l'antisimetria. Un bon ordre pot ser infinit i no tenir màxim.}

\pregunta El segment inicial \(W_{<w}\) és\ldots
\begin{itemize}
\opcio{a}{la part \(\{u\in W : u<w\}\).}
\opcio{b}{tot \(W\) tret de \(w\).}
\opcio{c}{el conjunt \(\{u : u>w\}\).}
\opcio{d}{el mínim de \(W\).}
\end{itemize}
\feedback
\justificacio{El segment és la part dels elements estrictament anteriors a \(w\), imatge inversa de \({<}\) al llarg de \(\langle\id,w\circ{!}\rangle\). No és el complement del punt.}

\pregunta En un bon ordre, les parts inicials pròpies\ldots
\begin{itemize}
\opcio{a}{no existeixen.}
\opcio{b}{poden no tenir mínim.}
\opcio{c}{són exactament els segments \(W_{<w}\), amb \(w=\min(W\setminus S)\).}
\opcio{d}{coincideixen amb tot \(W\).}
\end{itemize}
\feedback
\justificacio{Si \(S\) és inicial i pròpia, el mínim \(w\) del complement compleix \(S=W_{<w}\); i aquest \(w\) és únic perquè un segment determina el seu punt.}

\pregunta La \emph{rigidesa} d'un bon ordre diu que, per a \(f\) creixent i injectiva,\ldots
\begin{itemize}
\opcio{a}{\(f\) és exhaustiva.}
\opcio{b}{cap punt no baixa: \(w\le f(w)\) per a tot \(w\).}
\opcio{c}{\(f\) té un punt fix.}
\opcio{d}{\(f(w)<w\) per a algun \(w\).}
\end{itemize}
\feedback
\justificacio{Si la part dels punts que baixen fos no buida, el seu mínim \(w_0\) donaria \(f(w_0)<w_0\) i, aplicant \(f\), un punt que baixa encara més petit. Per tant cap punt no baixa.}

\pregunta Entre dos bons ordres, el nombre d'isomorfismes d'ordre és\ldots
\begin{itemize}
\opcio{a}{com a molt un.}
\opcio{b}{sempre exactament un.}
\opcio{c}{infinit.}
\opcio{d}{depèn de l'elecció.}
\end{itemize}
\feedback
\justificacio{Si \(f,g\) són isomorfismes d'ordre, \(g^{-1}\circ f\) i la seva inversa són automorfismes i cap punt no baixa en cap dels dos: la composició és la identitat i \(f=g\). Pot no haver-n'hi cap.}

\pregunta Cap bon ordre\ldots
\begin{itemize}
\opcio{a}{és total.}
\opcio{b}{té element mínim.}
\opcio{c}{és isomorf a cap segment inicial propi seu.}
\opcio{d}{admet aplicacions creixents.}
\end{itemize}
\feedback
\justificacio{Un tal isomorfisme donaria una creixent injectiva \(W\to W\) amb \(f(w)<w\), contra la rigidesa. És el que fa consistent comparar bons ordres per segments.}

\pregunta La comparació \enquote{ser un segment de} sobre bons ordres és\ldots
\begin{itemize}
\opcio{a}{només reflexiva.}
\opcio{b}{antisimètrica i transitiva (llevat d'isomorfisme d'ordre).}
\opcio{c}{simètrica.}
\opcio{d}{una relació d'equivalència.}
\end{itemize}
\feedback
\justificacio{Dues dominacions oposades componen un isomorfisme sobre un segment inicial de si mateix, que no pot ser propi; per tant és un automorfisme, la identitat. És un ordre llevat d'isomorfisme.}

\pregunta La comparabilitat total de dos bons ordres (que almenys una situació de la tricotomia es doni)\ldots
\begin{itemize}
\opcio{a}{és falsa en general.}
\opcio{b}{requereix l'axioma d'elecció.}
\opcio{c}{s'obté per recursió transfinita, sense elecció (el pas al mínim és canònic).}
\opcio{d}{és equivalent al teorema de Cantor.}
\end{itemize}
\feedback
\justificacio{L'aparellament es construeix per recursió prenent a cada pas el mínim encara lliure: una tria canònica, no una elecció. La comparabilitat de conjunts \emph{nus} sí que demana l'elecció.}

\pregunta La totalitat d'un ordre, \enquote{\(u=w\) o \(u<w\) o \(w<u\)}, s'enuncia per a punts globals. Què diu com a afirmació sobre parts?
\begin{itemize}
\opcio{a}{Que \({<}\) és tot \(W\times W\).}
\opcio{b}{Que \(\Delta\cup{<}\cup\operatorname{Inv}(<)\) és la part total de \(W\times W\); un element generalitzat pot quedar repartit entre els tres membres.}
\opcio{c}{Que \({<}\) és una relació d'equivalència.}
\opcio{d}{Que tot element generalitzat cau sencer en un dels tres membres.}
\end{itemize}
\feedback
\justificacio{Les condicions disjuntives són unions de parts, i la separació global permet comprovar-les sobre punts; però sobre un estadi general \enquote{o} no exclou, i (d) és fals. És la mateixa situació que a l'ordre de \(\N\).}

\pregunta Quina d'aquestes propietats \emph{no} val en un ordre total qualsevol, sinó només amb el bon ordre?
\begin{itemize}
\opcio{a}{Que tota aplicació creixent sigui injectiva.}
\opcio{b}{Que cap punt baixi sota una creixent injectiva (rigidesa) i, per tant, que hi hagi com a molt un isomorfisme d'ordre.}
\opcio{c}{Que l'ordre sigui transitiu.}
\opcio{d}{Que els segments inicials siguin parts.}
\end{itemize}
\feedback
\justificacio{La injectivitat de les creixents i la transitivitat valen en qualsevol ordre total; la rigidesa necessita el mínim d'una part (a \(\Z\) la translació \(n\mapsto n+1\) és creixent i injectiva i fa baixar tots els punts, en el sentit que cap punt no queda fix).}

\end{enumerate}
\clauestatica
\finalitzaTest
